% Edition 0.5: falsifiable research charter, not a GUT or ToE claim.
\chapter[Conditions for a GUT or ToE]{Beyond the Matrix Model: Conditions for a GUT, a Theory of Everything, and Genuine Physical Emergence}
\label{ch:unification-tests}

\section{A precise hierarchy of claims}
We distinguish four successive levels:
\begin{description}
\item[Mathematical structure.] An exact algebraic, group-theoretic,
functional-analytic or topological result within its stated carrier.
\item[Source-bound physical model.] A theorem about the actual
relative SU(2) quantum mechanics, with the original gauge action,
Hamiltonian and operator domains preserved.
\item[Emergent effective physics.] A derivation of lower-dimensional
fields and interactions from the source theory, together with
controlled limits and effective observables.
\item[Empirical unification.] A predictive candidate consistent with
the observed gauge symmetries, chiral matter, interactions and
experimental evidence, and, for a Theory of Everything, with a
credible quantum gravitational/spacetime sector as well.
\end{description}
Proving a result at one level does not automatically promote it to
the next. In particular, the original BFSS eleven-dimensional
M-theory conjecture was about the large-$N$ matrix model;
our fixed relative SU(2) construction does not, by adopting its
equations, prove that conjecture or establish a novel lower-dimensional
unification theory.

\section{The proposed bridge to four-dimensional particle physics}
A genuine GUT candidate must define a physical four-dimensional
spacetime or effective field-theory regime with Lorentz symmetry,
identify the internal gauge group and its breaking pattern, and
obtain the correct chiral fermion representations and hypercharges.
It must supply gauge anomaly cancellation, appropriate interactions,
and a clear account of why there are exactly three observed families.
A resemblance between a real spatial Spin(9) spinor and a complex
internal Spin(10) spinor does not discharge any of these requirements.

A more ambitious Theory of Everything requires a physically defined
spacetime/gravitational sector, quantum consistency, a principled
relation to known low-energy gravity, and a consistent coupling to
matter. The original fixed-$N$ relative matrix quantum mechanics is
mathematically well-defined in its existing Hilbert-space sense but
is not itself a derived local four-dimensional quantum field theory
or an established account of cosmological spacetime.

\section{Category theory as a validity test, not a substitute for dynamics}
The simultaneous source Spin(9) and SU(2) representations, the
original Clifford/Fock operator functoriality, and the closed physical
Hamiltonian suggest several candidate commuting diagrams. For
example, a prospective field $\mathcal O(X)$ should satisfy a
\emph{gauge-equivariance} identity of the form
\begin{equation}\label{eq:unification-gauge-natural}
 \mathcal O(q\!\cdot\!X)=U_F(q)\mathcal O(X)U_F(q)^{-1},
\end{equation}
where $U_F$ is the actual original Fock gauge unitary and the
left-hand source action is the genuine bosonic one. A prospective
spatial covariance theorem has a distinct source Spin(9) action.
Proving that \eqref{eq:unification-gauge-natural} holds for selected
Fock bilinears is a valuable mathematical result, but even exact
gauge covariance need not imply that $\mathcal O$ commutes with the
Hamiltonian or defines a new particle species.

The category-theoretic role is to expose missing naturality,
coherence and descent conditions. The negative Hopf-projector
example of Chapter~\ref{ch:unification-topology} is a model of
this discipline: it rejects a too-naive labelled subobject, while
permitting a different equivariant or stack-like object that can be
specified rigorously. No universal ``category of everything'' is
posited or admitted as an empirical framework merely because a
diagram can be drawn.

\section{Source-faithful research gates and kill criteria}
We organize the unification inquiry into five independent lanes:
\begin{enumerate}[leftmargin=*,itemsep=.5em]
\item \textbf{Representations.} Compare the source Spin(9)
$16_{\mathbb R}$ and its actual Spin(8) eigenspaces with abstract
Spin(10) complex chiral spinors. Reject any construction that
identifies spatial rotation with internal gauge symmetry without a
physical bridge.
\item \textbf{Physical families and charges.} Search for
three gauge-invariant chiral sectors with correctly normalized
charges, not three rotated color basis vectors. Record a no-go for
any source family-projector proposal that fails the original SU(2)
commutant test.
\item \textbf{Topology.} Build exact projectors, Hopf-type clutching
maps and gauge-groupoid descent objects. Reject independently labelled
$E_+$ sectors on the coarse projective quotient when they fail
naturality or characteristic-class descent.
\item \textbf{Dynamics.} For any proposed bundle, sector or composite,
construct its operator on the original physical space, establish
domain compatibility, and identify a conserved or dynamical role.
Reject a physical interpretation when only a static finite-dimensional
identity is available.
\item \textbf{Phenomenology.} Derive a falsifiable observable or
effective theory, compare it with known data, and invite independent
criticism. Do not promote a mathematical coincidence to a GUT or ToE
without prediction or physical realization.
\end{enumerate}

The original research on NFC, nested Hopf fibrations, protected
persistence, skyrmionic configurations and charge columns may provide
hypotheses for these lanes. No preferred hypothesis receives an
exception from the source, covariance, topology, domain or empirical
tests. A rejected correspondence can still be a useful theorem
about the precise conditions under which an identification is
impossible.

\section{Editorial independence from the eventual external monograph}
This internal text will preserve speculative motivations, failure
analyses, alternate representations and negative source checks.
The external manuscript should include an exploratory unification
claim only if it matures into a precisely stated and verified theorem
with appropriate physical interpretation. A mathematically complete
relative SU(2) spectral result would already be significant even if
none of the prospective GUT or ToE directions survives.

\formalrecord{Research-only sources:
\path{research/unification/BFSS_THREE_BY_SIXTEEN_SPIN9_SPIN10_FAMILY_FEASIBILITY_2026_11_10.md};
\path{research/unification/BFSS_ORIGINAL_MAJORANA_BILINEARS_RP8_TWISTED_QUADRATIC_BUNDLES_2026_11_10.md};
\path{research/unification/CATEGORY_THEORY_GAUGE_DESCENT_HOPF_NO_GO_2026_11_10.md};
\path{research/unification/TRUTH_FIRST_GUT_TOE_EVIDENCE_TRANSFER_AND_GLOBAL_GROUP_AUDIT_2026_11_10.md}.
No new strict Lean unification theorem, phenomenological prediction,
GUT or ToE is claimed.}
